\section{Análisis}

Prueba a 1\,kHz con \texttt{visor\_enlace}, con el FT2232H y la Curiosity
Nano en un solo hub USB y sin cable de GND entre las tarjetas. Con
\texttt{prueba\_enlace} en la misma conexión no hubo errores
(\texttt{dspic\_1khz} de esta carpeta); con el visor, que carga la PC y la
GPU, hubo 25 respuestas con CRC incorrecto (0.025\,\%), y en la respuesta
siguiente el dsPIC reportó longitud incorrecta para la misma transferencia.

Las respuestas dañadas muestran la causa. Por ejemplo, se esperaba
\texttt{5A AD 00 00 00 00 00} y se recibió \texttt{5A B0 05 AA D0 00 00 00}:
a partir del byte 1 aparece de nuevo el inicio de la respuesta
(\texttt{5A AD}) recorrido 4 bits. En otras el recorrido es de 3 bits. Sólo
la interrupción INT1 vuelve a armar la respuesta desde el principio, por lo
que un pulso falso en \texttt{CS} a mitad de la transferencia se toma como
fin de trama: la interrupción reinicia SPI1 y rearma el DMA mientras la PC
sigue transfiriendo. La corrupción empieza en el byte 1 en 13 casos, en el 2
en uno y en el byte 7 en 11, según el momento del pulso; es el mismo patrón
de \texttt{resultados/2026-10-05\_pll}.

Con el cable de GND la misma prueba no tuvo errores
(\texttt{resultados/2026-10-06\_gnd}). La bandera de vigilancia de la trama 1
viene de la inactividad previa a la prueba y no se cuenta.
